\chapter{Hypothesis Testing}
\label{ch:hypothesis-testing}

Chapters 7 and 8 built intervals: ranges of parameter values that are compatible
with the data. This chapter asks a sharper, yes-or-no style question. A friend
says her new checkout button is faster, or that she can predict coin flips, or that
a drug lowers blood pressure. The data show \emph{some} effect. Could plain luck
have produced it? A hypothesis test answers that by playing a game with a skeptic.

\section{The skeptic's game}
\label{sec:ht-skeptic}

Here is the puzzle. Your friend claims she can predict coin flips. You flip a fair
coin $20$ times; before each flip she calls heads or tails. She gets $16$ of the
$20$ calls right.

Is she gifted? You cannot answer that directly, because ``gifted'' is vague. Gifted
at the level of $60\%$? $75\%$? $99\%$? Each would predict a different number of
hits, and you have no way to pick among them. But there is one world you can describe
exactly: the world where she has no gift at all and each call is a fair
$50/50$ guess. Call it the \textbf{null world}, the boring world where nothing is
going on.

The skeptic's move is to say: ``Let us assume she is guessing, and see how
surprising $16$ would be.'' In the null world, the number of correct calls is a
count of heads in $20$ fair flips. You can simulate it. Imagine $1000$ people with
no gift at all, each making $20$ blind guesses against $20$ coin flips and writing
down how many they got right. Figure~\ref{fig:ht-null-world} is the histogram of
those $1000$ scores (a real simulation with a fixed seed).

\begin{figure}[H]
\centering
\begin{tikzpicture}
\begin{axis}[
  width=12cm, height=5.6cm, ybar, bar width=0.38cm,
  xlabel={correct calls out of 20}, ylabel={number of guessers (of 1000)},
  xmin=-0.5, xmax=20.5, ymin=0, ymax=235, xtick={0,2,...,20},
  axis lines=left, title={The null world: 1000 people guessing at random},
  title style={font=\small}, clip=false,
]
\addplot[fill=histfill, draw=distblue] coordinates {
(0,0)(1,0)(2,1)(3,2)(4,3)(5,18)(6,22)(7,71)(8,115)(9,151)(10,186)
(11,161)(12,141)(13,76)(14,30)(15,17)(16,5)(17,1)(18,0)(19,0)(20,0)};
\draw[decorate, decoration={brace, amplitude=5pt}, thick, structurepurple]
  (axis cs:7.6,195) -- (axis cs:12.4,195);
\node[structurepurple, font=\small, anchor=south] at (axis cs:10,203)
  {8 to 12 correct: 754 of the 1000 guessers};
\end{axis}
\end{tikzpicture}
\caption{What ``nothing is going on'' looks like. Three quarters of the random
guessers score between $8$ and $12$ out of $20$, and scores far from $10$ are rare.
This picture is the \emph{null distribution}; it was computed without knowing
anything about your friend.}
\label{fig:ht-null-world}
\end{figure}

Now drop her score onto the picture. In Figure~\ref{fig:ht-observed} the bars at
$16$ or more are shaded red. Only $5+1=6$ of the $1000$ random guessers did as well
as she did.

\begin{figure}[H]
\centering
\begin{tikzpicture}
\begin{axis}[
  name=main, width=12cm, height=5.6cm, ybar, bar width=0.38cm,
  xlabel={correct calls out of 20}, ylabel={number of guessers (of 1000)},
  xmin=-0.5, xmax=20.5, ymin=0, ymax=240, xtick={0,2,...,20},
  axis lines=left, clip=false,
]
\addplot[fill=histfill, draw=distblue, bar shift=0pt] coordinates {
(2,1)(3,2)(4,3)(5,18)(6,22)(7,71)(8,115)(9,151)(10,186)
(11,161)(12,141)(13,76)(14,30)(15,17)};
\addplot[fill=trapred, draw=trapred!70!black, bar shift=0pt] coordinates {(16,5)(17,1)(18,0)};
\draw[trapred, thick] (axis cs:15.5,-4) rectangle (axis cs:20.5,14);
\end{axis}
\begin{axis}[
  at={(main.north east)}, anchor=north east, xshift=-0.1cm, yshift=-0.1cm,
  width=5.0cm, height=3.6cm, ybar, bar width=0.26cm,
  xmin=12.5, xmax=20.5, ymin=0, ymax=105, xtick={13,...,20}, ytick=\empty,
  axis lines=left, clip=false, tick label style={font=\scriptsize},
  title={tail zoomed in (counts on bars)}, title style={font=\small},
  nodes near coords, nodes near coords style={font=\scriptsize},
]
\addplot[fill=histfill, draw=distblue, bar shift=0pt] coordinates {(13,76)(14,30)(15,17)};
\addplot[fill=trapred, draw=trapred!70!black, bar shift=0pt, nodes near coords style={font=\scriptsize, trapred}]
  coordinates {(16,5)(17,1)(18,0)(19,0)(20,0)};
\node[trapred, font=\scriptsize, anchor=south, align=center] at (axis cs:18,40)
  {her score 16 or more:\\$5+1=6$ of 1000};
\end{axis}
\end{tikzpicture}
\caption{The observed result dropped into the null world. In the full histogram the
tail at $16$ or more (red box) is too thin to see, so the inset zooms in on it: only
$6$ of the $1000$ simulated guessers did as well as your friend. That fraction, about
$0.006$, measures how surprised a skeptic should be.}
\label{fig:ht-observed}
\end{figure}

That fraction is the \textbf{p-value}: the share of the null world that produces a
result at least as extreme as the one you saw. The simulation gives $0.006$. The
exact binomial calculation (Probability From the Ground Up, Ch.~7) gives
\[
  \P(X \ge 16) = \sum_{k=16}^{20} \binom{20}{k} 0.5^{20}
  = \frac{6196}{1048576} \approx 0.0059 .
\]

Now the skeptic has two choices. Either your friend really is a random guesser and
a $1$-in-$170$ fluke just happened, or the null world is the wrong description of
her. Most people find the second option more believable. That is the whole logic of
a hypothesis test.

\begin{intuitionbox}
A hypothesis test is a proof by contradiction that is allowed to fail with small
probability. In a mathematical proof you assume a statement, derive something
impossible, and conclude the statement is false. Here you assume ``nothing is going
on'', derive that the data you saw would be \emph{very unlikely}, and conclude that
``nothing is going on'' is hard to believe. ``Impossible'' becomes ``improbable''.
\end{intuitionbox}

Why assume the boring hypothesis and not the exciting one? Because the boring one is
the only one specific enough to compute with. Try drawing the exciting world, ``she
has a gift'', and you get stuck at once: to draw a histogram you must pick a hit rate $\theta$ (theta), her
true long-run fraction of correct calls.
Figure~\ref{fig:ht-alt-family} draws three of the infinitely many choices.

\begin{figure}[H]
\centering
\begin{tikzpicture}
\begin{axis}[
  width=12cm, height=5.8cm, xlabel={correct calls out of 20}, ylabel={probability},
  xmin=-0.5, xmax=20.5, ymin=0, ymax=0.32, xtick={0,2,...,20}, axis lines=left,
  yticklabel style={/pgf/number format/fixed},
  legend style={font=\footnotesize, draw=none, fill=white, at={(0.02,0.98)}, anchor=north west},
  legend cell align=left,
]
\addplot[ybar, bar width=0.36cm, fill=histfill, draw=distblue, area legend] coordinates {
(2,0.0002)(3,0.0011)(4,0.0046)(5,0.0148)(6,0.0370)(7,0.0739)(8,0.1201)(9,0.1602)
(10,0.1762)(11,0.1602)(12,0.1201)(13,0.0739)(14,0.0370)(15,0.0148)(16,0.0046)(17,0.0011)(18,0.0002)};
\addlegendentry{null, $\theta=0.5$: one exact picture}
\addplot[thick, mark=*, mark size=1.3pt, targetgreen!80!black, dotted] coordinates {
(4,0.0003)(5,0.0013)(6,0.0049)(7,0.0146)(8,0.0355)(9,0.0710)(10,0.1171)(11,0.1597)
(12,0.1797)(13,0.1659)(14,0.1244)(15,0.0746)(16,0.0350)(17,0.0123)(18,0.0031)(19,0.0005)};
\addlegendentry{gift, $\theta=0.6$?}
\addplot[thick, mark=square*, mark size=1.3pt, targetgreen!80!black, dashed] coordinates {
(7,0.0002)(8,0.0008)(9,0.0030)(10,0.0099)(11,0.0271)(12,0.0609)(13,0.1124)(14,0.1686)
(15,0.2023)(16,0.1897)(17,0.1339)(18,0.0669)(19,0.0211)(20,0.0032)};
\addlegendentry{gift, $\theta=0.75$?}
\addplot[thick, mark=triangle*, mark size=1.6pt, targetgreen!80!black] coordinates {
(11,0.0001)(12,0.0004)(13,0.0020)(14,0.0089)(15,0.0319)(16,0.0898)(17,0.1901)(18,0.2852)
(19,0.2702)(20,0.1216)};
\addlegendentry{gift, $\theta=0.9$?}
\draw[trapred, thick, densely dashed] (axis cs:16,0) -- (axis cs:16,0.28);
\node[trapred, font=\small, anchor=south] at (axis cs:16,0.28) {her 16};
\end{axis}
\end{tikzpicture}
\caption{The null hypothesis is one world; ``she has a gift'' is a whole family of
worlds. Blue bars: guessing ($\theta=0.5$). Green curves: three of the infinitely many
gifted worlds, each a different prediction. Only the blue picture is pinned down well
enough to measure surprise against.}
\label{fig:ht-alt-family}
\end{figure}

``No gift'' pins $\theta$ to
exactly $0.5$ and produces one histogram. ``Some gift'' means $\theta > 0.5$, which
is every green curve at once and no histogram in particular. The test needs one
concrete world in which to measure surprise, and only the null hypothesis supplies
one. That is the asymmetry at the heart of testing: you can argue \emph{against} the
null world by showing your data look out of place in it, but there is no single
alternative world to argue \emph{for}.

A courtroom works the same way. The defendant is presumed innocent (the null).
The prosecution must show the evidence would be very unlikely if the defendant were
innocent. A verdict of ``not guilty'' does not mean ``proven innocent''; it means the
evidence was not strong enough. Hold on to that last point. It is the subject of
Section~\ref{sec:ht-absence}.

\subsection{Drawing a line in advance}

Someone has to decide how small is small enough. The convention is to fix a
threshold $\alpha$ (alpha), the \textbf{significance level}, before looking at the
data. $\alpha$ is the fraction of the null world you are willing to call
``surprising''. With $\alpha = 0.05$ and a one-sided test you reject the null when
your friend's score lands in the top $5\%$ of the null world.

For $20$ coin calls the counts are whole numbers, so the cut-off cannot hit $5\%$
exactly. $\P(X\ge 14) = 0.058$ is too big, and $\P(X\ge 15) = 0.021$ fits. So the
rule is: \emph{reject ``no gift'' if she gets $15$ or more right}.
Figure~\ref{fig:ht-threshold} shows this rejection region with exact probabilities.

\begin{figure}[H]
\centering
\begin{tikzpicture}
\begin{axis}[
  width=12cm, height=5.4cm, ybar, bar width=0.38cm,
  xlabel={correct calls out of 20}, ylabel={probability under the null},
  xmin=-0.5, xmax=20.5, ymin=0, ymax=0.21, xtick={0,2,...,20},
  axis lines=left, clip=false, yticklabel style={/pgf/number format/fixed},
]
\addplot[fill=histfill, draw=distblue, bar shift=0pt] coordinates {
(2,0.0002)(3,0.0011)(4,0.0046)(5,0.0148)(6,0.0370)(7,0.0739)(8,0.1201)(9,0.1602)
(10,0.1762)(11,0.1602)(12,0.1201)(13,0.0739)(14,0.0370)};
\addplot[fill=trapred, draw=trapred!70!black, bar shift=0pt] coordinates {
(15,0.0148)(16,0.0046)(17,0.0011)(18,0.0002)};
\draw[densityred, dashed, thick] (axis cs:14.5,0) -- (axis cs:14.5,0.15);
\node[densityred, anchor=south west, font=\small, align=left] at (axis cs:14.6,0.075)
  {rejection region\\$X\ge 15$\\total $0.021\le\alpha$};
\end{axis}
\end{tikzpicture}
\caption{Fixing $\alpha=0.05$ before the experiment turns the null world into a
decision rule. Results in the red region lead to rejecting the null. If the null is
true, this happens with probability $0.021$, never more than $\alpha$.}
\label{fig:ht-threshold}
\end{figure}

The whole procedure fits in one storyboard (Figure~\ref{fig:ht-flow}). Every test in
this chapter follows panels 1--5; only the statistic and the null world change.
Panel 6 is the step the null world cannot do alone: to ask how often the test would
catch a \emph{real} gift, you must pick one gifted world and overlay it. That
overlay is called power, and Section~\ref{sec:ht-power} develops it.

\begin{figure}[H]
\centering
\begin{tikzpicture}[sbcap/.style={font=\footnotesize, align=center, text width=3.7cm, anchor=north}]
\begin{scope}
\draw[rounded corners, fill=sourceblue!8, draw=sourceblue] (0.05,0.1) rectangle (3.55,2.1);
\node[font=\small, align=center] at (1.8,1.55) {Skeptic: ``she is\\just guessing''};
\foreach \i/\s in {0/H,1/T,2/T,3/H,4/T} {
  \draw[fill=quotientorange!25, draw=quotientorange!80!black] (0.75+\i*0.52,0.55) circle (0.2);
  \node[font=\scriptsize] at (0.75+\i*0.52,0.55) {\s};
}
\node[sbcap] at (1.8,-0.4) {1. Assume nothing is going on: $H_0\colon\theta=0.5$};
\end{scope}
\begin{scope}[xshift=4.2cm]
\draw (-0.1,0) -- (3.6,0);
\foreach \k/\p in {2/0.0002,3/0.0011,4/0.0046,5/0.0148,6/0.0370,7/0.0739,8/0.1201,9/0.1602,10/0.1762,11/0.1602,12/0.1201,13/0.0739,14/0.0370,15/0.0148,16/0.0046,17/0.0011,18/0.0002}
  \filldraw[fill=histfill, draw=distblue, line width=0.3pt] ({(\k-2)*0.19-0.075},0) rectangle ({(\k-2)*0.19+0.075},{\p*10});
\node[font=\scriptsize, anchor=north] at (1.52,0) {10};
\node[sbcap] at (1.8,-0.4) {2. Picture the null world: scores of pure guessers};
\end{scope}
\begin{scope}[xshift=8.4cm]
\draw (-0.1,0) -- (3.6,0);
\foreach \k/\p in {2/0.0002,3/0.0011,4/0.0046,5/0.0148,6/0.0370,7/0.0739,8/0.1201,9/0.1602,10/0.1762,11/0.1602,12/0.1201,13/0.0739,14/0.0370,15/0.0148,16/0.0046,17/0.0011,18/0.0002}
  \filldraw[fill=histfill, draw=distblue, line width=0.3pt] ({(\k-2)*0.19-0.075},0) rectangle ({(\k-2)*0.19+0.075},{\p*10});
\draw[trapred, very thick, -{Stealth}] (2.66,1.45) -- (2.66,0.08);
\node[trapred, font=\scriptsize, anchor=south] at (2.66,1.45) {her 16};
\node[font=\scriptsize, anchor=north] at (1.52,0) {10};
\node[sbcap] at (1.8,-0.4) {3. Drop the observed result onto it};
\end{scope}
\begin{scope}[yshift=-3.9cm]
\draw (-0.1,0) -- (3.6,0);
\foreach \k/\p in {2/0.0002,3/0.0011,4/0.0046,5/0.0148,6/0.0370,7/0.0739,8/0.1201,9/0.1602,10/0.1762,11/0.1602,12/0.1201,13/0.0739,14/0.0370,15/0.0148}
  \filldraw[fill=histfill, draw=distblue, line width=0.3pt] ({(\k-2)*0.19-0.075},0) rectangle ({(\k-2)*0.19+0.075},{\p*10});
\foreach \k/\p in {16/0.0046,17/0.0011,18/0.0002}
  \filldraw[fill=trapred, draw=trapred!70!black, line width=0.3pt] ({(\k-2)*0.19-0.075},0) rectangle ({(\k-2)*0.19+0.075},{\p*10});
\draw[trapred, line width=2.5pt] (2.565,0) -- (3.515,0);
\node[trapred, font=\scriptsize, align=center, anchor=south] at (2.9,0.55) {tail $\ge 16$\\$p=0.006$};
\draw[trapred] (2.9,0.55) -- (2.9,0.06);
\node[font=\scriptsize, anchor=north] at (1.52,0) {10};
\node[sbcap] at (1.8,-0.4) {4. Shade the tail at least as extreme: its area is the p-value};
\end{scope}
\begin{scope}[xshift=4.2cm, yshift=-3.9cm]
\draw (-0.1,0) -- (3.6,0);
\foreach \k/\p in {2/0.0002,3/0.0011,4/0.0046,5/0.0148,6/0.0370,7/0.0739,8/0.1201,9/0.1602,10/0.1762,11/0.1602,12/0.1201,13/0.0739,14/0.0370}
  \filldraw[fill=histfill, draw=distblue, line width=0.3pt] ({(\k-2)*0.19-0.075},0) rectangle ({(\k-2)*0.19+0.075},{\p*10});
\foreach \k/\p in {15/0.0148,16/0.0046,17/0.0011,18/0.0002}
  \filldraw[fill=trapred, draw=trapred!70!black, line width=0.3pt] ({(\k-2)*0.19-0.075},0) rectangle ({(\k-2)*0.19+0.075},{\p*10});
\draw[densityred, dashed, thick] (2.375,0) -- (2.375,1.75);
\draw[trapred, line width=2.5pt] (2.39,0) -- (3.515,0);
\node[densityred, font=\scriptsize, anchor=west, align=left] at (2.45,1.2) {reject\\if $\ge 15$};
\node[font=\scriptsize, anchor=north] at (1.52,0) {10};
\node[sbcap] at (1.8,-0.4) {5. Fix the $\alpha$ line in advance; reject beyond it};
\end{scope}
\begin{scope}[xshift=8.4cm, yshift=-3.9cm]
\draw (-0.1,0) -- (3.6,0);
\foreach \k/\p in {2/0.0002,3/0.0011,4/0.0046,5/0.0148,6/0.0370,7/0.0739,8/0.1201,9/0.1602,10/0.1762,11/0.1602,12/0.1201,13/0.0739,14/0.0370,15/0.0148,16/0.0046,17/0.0011,18/0.0002}
  \filldraw[fill=histfill, draw=distblue, line width=0.3pt] ({(\k-2)*0.19-0.09},0) rectangle ({(\k-2)*0.19},{\p*10});
\foreach \k/\p in {7/0.0002,8/0.0008,9/0.0030,10/0.0099,11/0.0271,12/0.0609,13/0.1124,14/0.1686}
  \filldraw[fill=bootorange!50, draw=bootorange, line width=0.3pt] ({(\k-2)*0.19},0) rectangle ({(\k-2)*0.19+0.09},{\p*10});
\foreach \k/\p in {15/0.2023,16/0.1897,17/0.1339,18/0.0669,19/0.0211,20/0.0032}
  \filldraw[fill=targetgreen!60, draw=targetgreen!60!black, line width=0.3pt] ({(\k-2)*0.19},0) rectangle ({(\k-2)*0.19+0.09},{\p*10});
\draw[densityred, dashed, thick] (2.375,0) -- (2.375,2.15);
\node[targetgreen!50!black, font=\scriptsize, anchor=west] at (2.45,2.25) {power $0.62$};
\node[distblue, font=\scriptsize, anchor=east] at (1.25,1.8) {guessing};
\node[font=\scriptsize, anchor=north] at (1.52,0) {10};
\node[sbcap] at (1.8,-0.4) {6. Overlay a gifted world ($\theta=0.75$): green share $=$ power};
\end{scope}
\end{tikzpicture}
\caption{The logic of every hypothesis test as a storyboard. Panels 1--5 use only
the null world: assume it, draw it, place the data, measure the tail, compare with a
line drawn in advance. Panel 6 overlays one gifted world (orange and green bars,
next to the blue guessing bars) to see how often a real effect would land past the
line; the green share is the power.}
\label{fig:ht-flow}
\end{figure}

\begin{recoverybox}
Pause and think: if your friend had scored $13$ of $20$, what would the skeptic
conclude? (Look at Figure~\ref{fig:ht-threshold}: $13$ sits in the blue bulk.
$\P(X\ge 13)\approx 0.13$, so the data are unsurprising under the null and you do not
reject. You also have not shown she lacks a gift.)
\end{recoverybox}

\section{Absence of evidence is not evidence of absence}
\label{sec:ht-absence}

Suppose a second friend makes $10$ calls and gets $6$ right, a $60\%$ hit rate. A
third friend makes $1000$ calls and gets $600$ right, also $60\%$. Same rate. Very
different verdicts.

\begin{figure}[H]
\centering
\begin{tikzpicture}
\begin{axis}[
  name=left, width=6.3cm, height=5cm, ybar, bar width=0.17cm,
  xlabel={fraction correct}, ylabel={probability}, title={$n=10$: 6 right},
  title style={font=\small}, xmin=-0.05, xmax=1.05, ymin=0, ymax=0.36,
  axis lines=left, xtick={0,0.5,1}, ytick={0,0.1,0.2,0.3},
  yticklabel style={/pgf/number format/fixed}, clip=false,
]
\addplot[fill=histfill, draw=distblue, bar shift=0pt] coordinates {
(0,0.0010)(0.1,0.0098)(0.2,0.0439)(0.3,0.1172)(0.4,0.2051)(0.5,0.2461)};
\addplot[fill=trapred!60, draw=trapred, bar shift=0pt] coordinates {
(0.6,0.2051)(0.7,0.1172)(0.8,0.0439)(0.9,0.0098)(1.0,0.0010)};
\draw[trapred, thick] (axis cs:0.6,0) -- (axis cs:0.6,0.27);
\node[trapred, font=\footnotesize, anchor=south, align=center] at (axis cs:0.6,0.27)
  {observed 0.6\\red tail: $p=0.38$};
\end{axis}
\begin{axis}[
  at={(left.east)}, anchor=west, xshift=1.3cm,
  width=6.3cm, height=5cm, xlabel={fraction correct}, ylabel={density},
  title={$n=1000$: 600 right}, title style={font=\small},
  xmin=0.4, xmax=0.65, ymin=0, ymax=33, axis lines=left, domain=0.42:0.58,
  samples=120, xtick={0.4,0.5,0.6}, clip=false,
  xticklabel style={/pgf/number format/fixed, /pgf/number format/precision=1},
]
\addplot[fill=histfill, draw=distblue, thick] {25.231*exp(-(x-0.5)^2/(2*0.00025))} \closedcycle;
\draw[trapred, thick] (axis cs:0.6,0) -- (axis cs:0.6,22);
\node[trapred, font=\footnotesize, anchor=south, align=center] at (axis cs:0.6,22)
  {observed 0.6\\$p\approx 10^{-10}$};
\draw[{Stealth}-{Stealth}, structurepurple] (axis cs:0.5,3) -- (axis cs:0.6,3);
\node[structurepurple, font=\scriptsize, anchor=south] at (axis cs:0.567,3) {6.3 SD};
\end{axis}
\end{tikzpicture}
\caption{Same observed hit rate ($60\%$), two sample sizes; each panel is the null
world for that many calls. With $10$ calls (left) the null world is wide and $0.6$
sits inside it: $38\%$ of pure guessers do at least as well. With $1000$ calls
(right) the null world is narrow and $0.6$ is $6.3$ standard deviations outside it.
Failing to reject on the left says ``too little data'', not ``no gift''.}
\label{fig:ht-absence}
\end{figure}

On the left of Figure~\ref{fig:ht-absence}, the null world for $n=10$ is so wide
that $6$ of $10$ is utterly ordinary: $\P(X \ge 6) = 0.377$. On the right, the null
world for $n=1000$ has standard deviation $\sqrt{0.25/1000}\approx 0.016$ on the
fraction scale, and $0.6$ sits more than six standard deviations out.

The second friend might have exactly the same skill as the third. Ten calls
cannot tell a $60\%$ predictor from a $50\%$ guesser. To see this, run the test not
only for the hit rate $0.5$ but for every candidate hit rate $\theta_0$, and keep
the ones the data fail to reject (two-sided at $\alpha=0.05$, exact binomial
probabilities, $2.5\%$ allowed in each tail). Figure~\ref{fig:ht-survivors} shows the survivors.

\begin{figure}[H]
\centering
\begin{tikzpicture}[x=10cm]
\foreach \y/\lo/\hi/\obs/\lab in {1.5/0.262/0.878/0.6/{6 of 10}, 0/0.569/0.631/0.6/{600 of 1000}} {
  \draw[->] (0,\y) -- (1.06,\y) node[right, font=\small] {$\theta_0$};
  \foreach \t in {0,0.25,0.5,0.75,1} \draw (\t,\y+0.08) -- (\t,\y-0.08) node[below, font=\footnotesize] {\t};
  \fill[targetgreen!45] (\lo,\y-0.12) rectangle (\hi,\y+0.12);
  \draw[targetgreen!60!black] (\lo,\y-0.12) rectangle (\hi,\y+0.12);
  \node[font=\small, anchor=east] at (-0.02,\y) {\lab};
}
\draw[trapred, thick, dashed] (0.5,-0.08) -- (0.5,0.35);
\draw[trapred, thick, dashed] (0.5,1.42) -- (0.5,2.25);
\node[trapred, font=\footnotesize, anchor=south] at (0.5,2.25) {guessing, $\theta_0=0.5$};
\node[targetgreen!50!black, font=\footnotesize, anchor=south west] at (0.56,1.65) {not rejected: $0.26$ to $0.88$};
\node[targetgreen!50!black, font=\footnotesize, anchor=south west] at (0.64,0.15) {not rejected: $0.57$ to $0.63$};
\end{tikzpicture}
\caption{Which hit rates survive the test? With $6$ of $10$ (top), every value from
$0.26$ to $0.88$ survives, including $0.5$ and also a strong gift. Failing to reject
$0.5$ cannot single it out. With $600$ of $1000$ (bottom) the survivors shrink to a
narrow band that excludes $0.5$.}
\label{fig:ht-survivors}
\end{figure}

With $6$ of $10$, guessing survives, but so does a $60\%$ gift and so does an $85\%$
gift. The data have not picked out $0.5$; they have failed to rule out a wide range
of worlds, and $0.5$ is one of them. Declaring ``$\theta=0.5$ is accepted'' would be
as arbitrary as declaring ``$\theta=0.85$ is accepted''. That is why a test can
\emph{reject} the null but never \emph{accept} it. When you fail to reject, the
honest report is ``the data are compatible with the null'', and also with
everything else in the green band. (The green band is a confidence interval; Section~\ref{sec:ht-duality}
shows that this is no coincidence.) Section~\ref{sec:ht-power} turns ``too little
data'' into a number called power.

\begin{trapbox}
``We found no significant difference, so the two versions perform the same.'' This is
the most common misreading of a test. A large p-value from a small experiment is
an absence of evidence. To claim two things are \emph{equal}, look at the
confidence interval for the difference (Chapter 7): if it is narrow and sits inside a
range you consider negligible, then you have shown equivalence.
\end{trapbox}

\section{A permutation test: the null world by shuffling}
\label{sec:ht-permutation}

Coin flips have a known null world. Most real questions do not. Here is a
technique that builds the null world out of the data itself.

An engineer rewrites the rendering path of a checkout page. She loads the page $12$
times: $6$ with the old code, $6$ with the new, in random order. Render times in ms:
\begin{center}
\begin{tabular}{lcccccc|c}
\toprule
old & 212 & 198 & 225 & 240 & 205 & 231 & mean $218.5$\\
new & 190 & 184 & 207 & 195 & 178 & 201 & mean $192.5$\\
\bottomrule
\end{tabular}
\end{center}
The new code is $26.0$ ms faster on average. Could the labels ``old'' and ``new'' be
irrelevant, with the $26$ ms gap an accident of which loads happened to land in which
group?

The null hypothesis here is: \emph{the code version makes no difference to render
time}. If that is true, each of the $12$ numbers would have come out the same no
matter which label it got. The labels are then arbitrary tags, and any other way of
assigning $6$ ``old'' and $6$ ``new'' tags was equally likely. So shuffle the labels
and recompute the difference (Figure~\ref{fig:ht-shuffle}).

\begin{figure}[H]
\centering
\begin{tikzpicture}[cell/.style={draw, minimum width=0.74cm, minimum height=0.55cm, font=\small, inner sep=1pt}]
\node[anchor=west, font=\small] at (-0.4,0.6) {Real labels: $\bar o - \bar n = 218.5-192.5 = 26.0$};
\foreach \v/\c [count=\i from 0] in {212/sourceblue,198/sourceblue,225/sourceblue,240/sourceblue,205/sourceblue,231/sourceblue,190/quotientorange,184/quotientorange,207/quotientorange,195/quotientorange,178/quotientorange,201/quotientorange}
  \node[cell, fill=\c!40] at (\i*0.86,0) {\v};
\node[anchor=west, font=\small] at (-0.4,-0.9) {Shuffle 1: $202.8-208.2 = -5.3$};
\foreach \v/\c [count=\i from 0] in {212/sourceblue,198/quotientorange,225/sourceblue,240/quotientorange,205/sourceblue,231/quotientorange,190/sourceblue,184/sourceblue,207/quotientorange,195/quotientorange,178/quotientorange,201/sourceblue}
  \node[cell, fill=\c!40] at (\i*0.86,-1.5) {\v};
\node[anchor=west, font=\small] at (-0.4,-2.4) {Shuffle 2: $206.3-204.7 = 1.7$};
\foreach \v/\c [count=\i from 0] in {212/quotientorange,198/sourceblue,225/quotientorange,240/sourceblue,205/quotientorange,231/sourceblue,190/quotientorange,184/sourceblue,207/sourceblue,195/quotientorange,178/sourceblue,201/quotientorange}
  \node[cell, fill=\c!40] at (\i*0.86,-3.0) {\v};
\node[cell, fill=sourceblue!40] at (0,-4.0) {};
\node[anchor=west, font=\small] at (0.45,-4.0) {labelled ``old''};
\node[cell, fill=quotientorange!45] at (3.6,-4.0) {};
\node[anchor=west, font=\small] at (4.05,-4.0) {labelled ``new''};
\end{tikzpicture}
\caption{Shuffling labels. The twelve measurements stay put; only the colours
(labels) move. Under the null the labels carry no information, so each shuffle is a
draw from the null world. Shuffled differences are typically a few ms, far from $26$.}
\label{fig:ht-shuffle}
\end{figure}

There are $\binom{12}{6} = 924$ ways to choose which $6$ of the $12$ values get the
``old'' label. With so few, we can compute the difference for every one of them.
Figure~\ref{fig:ht-perm-null} is the histogram of all $924$.

\begin{figure}[H]
\centering
\begin{tikzpicture}
\begin{axis}[
  width=12cm, height=5.6cm, ybar interval, xlabel={difference in means under shuffled labels (ms)},
  ylabel={number of labellings}, ymin=0, ymax=175, xmin=-32, xmax=32,
  xtick={-26,-13,0,13,26}, axis lines=left, clip=false,
  x tick label as interval=false, xmajorgrids=false,
]
\addplot[fill=histfill, draw=distblue] coordinates {
(-30,5)(-26,15)(-22,28)(-18,52)(-14,78)(-10,105)(-6,120)(-2,118)
(2,120)(6,105)(10,78)(14,52)(18,28)(22,15)(26,5)(30,5)};
\filldraw[fill=trapred, draw=trapred!70!black] (axis cs:-30,0) rectangle (axis cs:-26,5);
\filldraw[fill=trapred, draw=trapred!70!black] (axis cs:26,0) rectangle (axis cs:30,5);
\draw[trapred, very thick] (axis cs:26,0) -- (axis cs:26,128);
\draw[trapred, very thick, dashed] (axis cs:-26,0) -- (axis cs:-26,128);
\node[trapred, anchor=south east, font=\small, align=right] at (axis cs:27,128) {observed $26.0$\\5 labellings $\ge 26$};
\node[trapred, anchor=south west, font=\small, align=left] at (axis cs:-27,128) {mirror $-26.0$\\5 labellings $\le -26$};
\node[trapred, font=\small, anchor=south] at (axis cs:28,5) {5};
\node[trapred, font=\small, anchor=south] at (axis cs:-28,5) {5};
\end{axis}
\end{tikzpicture}
\caption{The exact null distribution of the difference in means, from all $924$
relabellings (bins $4$ ms wide, edges at $\pm 26$). The real difference, $26.0$ ms,
sits in the far right tail: the red bar holds the only $5$ labellings with a
difference of $26$ or more, and its mirror holds $5$ more with $-26$ or less.
The histogram is symmetric because swapping every label flips the sign.}
\label{fig:ht-perm-null}
\end{figure}

Count the tail. Exactly $5$ of the $924$ labellings give a difference of $26.0$ or
more, and $5$ more give $-26.0$ or less. So:
\[
  p_{\text{one-sided}} = \frac{5}{924} \approx 0.0054,
  \qquad
  p_{\text{two-sided}} = \frac{10}{924} \approx 0.011 .
\]
If the code change did nothing, a gap this large would show up in about $1$ random
labelling in $90$. That is strong evidence the new code is faster.

Notice what we did not need: no Normal distribution, no formula for a standard
error, no assumption about the shape of render times. We only needed one fact: under
the null, the labels are exchangeable (any relabelling is as likely as any other).
That fact holds because the engineer randomized which loads used which code
(Chapter 5).

\begin{lstlisting}[language=Python]
import numpy as np
old = np.array([212,198,225,240,205,231]); new = np.array([190,184,207,195,178,201])
pooled = np.concatenate([old, new]); obs = old.mean() - new.mean()
rng = np.random.default_rng(0); diffs = []
for _ in range(100_000):                     # random shuffles instead of all 924
    p = rng.permutation(pooled)
    diffs.append(p[:6].mean() - p[6:].mean())
print(np.mean(np.abs(diffs) >= obs))         # 0.011, two-sided p-value
\end{lstlisting}

With larger groups, enumerating every labelling is impossible
($\binom{100}{50}\approx 10^{29}$), so you shuffle a large random sample of
labellings instead, as in the code. The answer converges to the exact one.

\section{The vocabulary, precisely}
\label{sec:ht-vocab}

Every test has the same parts. Here they are, each tied to the examples above.

\begin{definition}[Hypotheses]
The \textbf{null hypothesis} $H_0$ is a statement specific enough to determine the
distribution of the data (or at least of a chosen statistic). The \textbf{alternative
hypothesis} $H_1$ is what you conclude if you reject $H_0$.
\end{definition}

For the coin: $H_0\colon \theta = 0.5$ versus $H_1\colon \theta > 0.5$, where
$\theta$ is her true hit rate. For the render times: $H_0\colon$ the code makes no
difference; in mean-based tests this becomes $H_0\colon \mu_{\text{old}} =
\mu_{\text{new}}$, with $H_1\colon \mu_{\text{old}} \ne \mu_{\text{new}}$.

\begin{definition}[Test statistic and null distribution]
A \textbf{test statistic} $T$ is a number computed from the data that tends to be
larger (or further from zero) when $H_1$ is true. Its distribution assuming $H_0$ is
the \textbf{null distribution}.
\end{definition}

The count of correct calls and the difference in group means are test statistics.
Figures~\ref{fig:ht-null-world} and~\ref{fig:ht-perm-null} are null distributions.

\begin{definition}[p-value]
Let $t_{\text{obs}}$ be the observed value of $T$. The \textbf{p-value} is
\[
  p = \P_{H_0}\bigl(T \text{ is at least as extreme as } t_{\text{obs}}\bigr),
\]
the probability, computed in the null world, of a result at least as extreme as the
one observed. ``Extreme'' means ``in the direction of $H_1$''.
\end{definition}

The notation $\P_{H_0}$ reads ``probability computed as if $H_0$ were true''. It is
a reminder that the p-value is a statement about the null world, not about the
real one.

\subsection{One-sided and two-sided}

``Extreme'' depends on $H_1$. If only an increase counts as evidence (the coin
friend claims to be \emph{better} than chance), extreme means large, and the p-value
is the area of the right tail. If either direction counts (the new code might be
slower or faster), extreme means far from the null value on either side, and the
p-value adds both tails.

\begin{figure}[H]
\centering
\begin{tikzpicture}
\begin{axis}[
  name=one, width=6.4cm, height=4.6cm, domain=-3.6:3.6, samples=100,
  axis lines=left, ymin=0, ymax=0.48, ytick=\empty, xtick={0,1.645},
  xticklabels={0,$1.645$}, title={One-sided, $\alpha=0.05$}, title style={font=\small},
  xlabel={$z$},
]
\addplot[fill=trapred!60, draw=none, domain=1.645:3.6] {0.39894*exp(-0.5*x^2)} \closedcycle;
\addplot[thick, distblue] {0.39894*exp(-0.5*x^2)};
\node[font=\footnotesize, trapred, anchor=south, align=center] at (axis cs:2.6,0.08) {area\\$0.05$};
\end{axis}
\begin{axis}[
  at={(one.east)}, anchor=west, xshift=0.9cm,
  width=6.4cm, height=4.6cm, domain=-3.6:3.6, samples=100,
  axis lines=left, ymin=0, ymax=0.48, ytick=\empty, xtick={-1.96,0,1.96},
  xticklabels={$-1.96$,0,$1.96$}, title={Two-sided, $\alpha=0.05$}, title style={font=\small},
  xlabel={$z$},
]
\addplot[fill=trapred!60, draw=none, domain=1.96:3.6] {0.39894*exp(-0.5*x^2)} \closedcycle;
\addplot[fill=trapred!60, draw=none, domain=-3.6:-1.96] {0.39894*exp(-0.5*x^2)} \closedcycle;
\addplot[thick, distblue] {0.39894*exp(-0.5*x^2)};
\node[font=\footnotesize, trapred, anchor=south, align=center] at (axis cs:2.75,0.07) {area\\$0.025$};
\node[font=\footnotesize, trapred, anchor=south, align=center] at (axis cs:-2.75,0.07) {area\\$0.025$};
\end{axis}
\end{tikzpicture}
\caption{Rejection regions (red) for a statistic that is $\Normal(0,1)$ under $H_0$;
both panels shade a total area of $0.05$. A one-sided test spends all of $\alpha$ in
one tail, so its cut-off ($1.645$) is closer to zero. A two-sided test splits
$\alpha$ between both tails, so each cut-off moves out to $\pm 1.96$.}
\label{fig:ht-sides}
\end{figure}

\begin{trapbox}
Choose the sides \emph{before} seeing the data. Running a two-sided test, getting
$p=0.08$, then switching to one-sided to get $p=0.04$ doubles your real false-alarm
rate. The default in practice is two-sided. Use one-sided only when an effect in the
other direction would lead to exactly the same action as no effect.
\end{trapbox}

\section{The z-test}
\label{sec:ht-z}

The permutation test built the null world by brute force. When the statistic is a
sample mean, Chapter 6 already told us what the null world looks like: by the
Central Limit Theorem, $\Xbar$ is approximately Normal with mean $\mu$ and standard
error $\sigma/\sqrt{n}$. That gives a null world from a formula.

\begin{example}[API latency]
Historically, an API's response time has mean $200$ ms and standard deviation
$\sigma = 40$ ms. After a dependency upgrade, a random sample of $n=64$ requests has
mean $\xbar = 211$ ms. Did the mean change?

\textbf{Hypotheses.} $H_0\colon \mu = 200$ versus $H_1\colon \mu \ne 200$.

\textbf{Null world.} If $H_0$ holds, $\Xbar \approx \Normal(200, \SE^2)$ with
$\SE = 40/\sqrt{64} = 5$ ms.

\textbf{Statistic.} Measure the observed mean in units of standard error:
\[
  z = \frac{\xbar - \mu_0}{\sigma/\sqrt n} = \frac{211-200}{5} = 2.2 .
\]
\textbf{p-value.} Under $H_0$, $Z\sim\Normal(0,1)$, so
\[
  p = \P(|Z| \ge 2.2) = 2\bigl(1-\Phi(2.2)\bigr) = 2(0.0139) = 0.028,
\]
where $\Phi$ is the standard Normal CDF (the area to the left of a point under the
bell curve). Since $0.028 < 0.05$, reject $H_0$ at the $5\%$ level: the
mean latency appears to have risen.
\end{example}

\begin{figure}[H]
\centering
\begin{tikzpicture}
\begin{axis}[
  width=12cm, height=5cm, domain=-4:4, samples=140, axis lines=left,
  ymin=0, ymax=0.45, xlabel={$z$ (and the corresponding $\bar x$ in ms)},
  ylabel={density under $H_0$}, xtick={-2.2,0,2.2},
  xticklabels={$-2.2$ \small(189),{0 \small(200)},$2.2$ \small(211)}, clip=false,
]
\addplot[fill=trapred!60, draw=none, domain=2.2:4] {0.39894*exp(-0.5*x^2)} \closedcycle;
\addplot[fill=trapred!60, draw=none, domain=-4:-2.2] {0.39894*exp(-0.5*x^2)} \closedcycle;
\addplot[very thick, distblue] {0.39894*exp(-0.5*x^2)};
\draw[trapred, thick] (axis cs:2.2,0) -- (axis cs:2.2,0.2);
\node[trapred, anchor=south west, font=\small, align=left] at (axis cs:2.3,0.06) {area\\$0.0139$};
\node[trapred, anchor=south east, font=\small, align=right] at (axis cs:-2.3,0.06) {area\\$0.0139$};
\node[trapred, anchor=south, font=\small] at (axis cs:2.2,0.2) {observed $z=2.2$};
\end{axis}
\end{tikzpicture}
\caption{The two-sided p-value for $z=2.2$ is the total red area, $0.028$. The left
tail counts because a drop to $189$ ms would have been equally surprising under
$H_0\colon\mu=200$.}
\label{fig:ht-ztail}
\end{figure}

The z-test needs $\sigma$ known. That is rare: here it came from a long history of
the same service. Usually $\sigma$ must be estimated from the same sample, which
changes the null world slightly. That is the t-test.

\section{The one-sample t-test}
\label{sec:ht-t1}

Replace the unknown $\sigma$ by the sample standard deviation $s$ (Chapter 2):
\[
  t = \frac{\xbar - \mu_0}{s/\sqrt n}.
\]
This looks like $z$, but its null distribution is wider. The denominator $s$ is
itself random: in some samples it comes out too small, which inflates $t$. The extra
spread is captured by the \textbf{$t$ distribution with $n-1$ degrees of freedom},
written $t_{n-1}$. ``Degrees of freedom'' is the $n-1$ from the variance formula: the
number of independent deviations left after estimating the mean.

\begin{figure}[H]
\centering
\begin{tikzpicture}
\begin{axis}[
  width=11cm, height=5cm, domain=-4.5:4.5, samples=150, axis lines=left,
  ymin=0, ymax=0.43, xlabel={$t$}, ylabel={density},
  legend style={font=\small, draw=none, fill=none}, legend pos=north west,
  legend cell align=left, clip=false,
]
\addplot[very thick, distblue] {0.39894*exp(-0.5*x^2)};
\addlegendentry{$\Normal(0,1)$}
\addplot[thick, quotientorange] {0.38803*(1+x^2/9)^(-5)};
\addlegendentry{$t_9$}
\addplot[thick, trapred, dashed] {0.35355*(1+x^2/2)^(-1.5)};
\addlegendentry{$t_2$}
\draw[distblue, thick] (axis cs:1.96,0) -- (axis cs:1.96,0.12);
\draw[quotientorange, thick] (axis cs:2.262,0) -- (axis cs:2.262,0.12);
\draw[distblue] (axis cs:1.96,0.12) -- (axis cs:2.75,0.22);
\node[distblue, font=\footnotesize, anchor=west] at (axis cs:2.75,0.22) {Normal 5\% cut-off: $1.96$};
\draw[quotientorange!80!black] (axis cs:2.262,0.12) -- (axis cs:2.75,0.14);
\node[quotientorange!80!black, font=\footnotesize, anchor=west] at (axis cs:2.75,0.14) {$t_9$ 5\% cut-off: $2.262$};
\end{axis}
\end{tikzpicture}
\caption{The $t$ distribution is a bell with heavier tails. With few degrees of
freedom (small samples) extreme values are more common, so the two-sided cut-off for
$\alpha=0.05$ (vertical lines) is further out: $2.262$ for $t_9$ versus $1.96$ for
the Normal. As $n$ grows, $t_{n-1}$ approaches $\Normal(0,1)$.}
\label{fig:ht-tdist}
\end{figure}

\begin{whybox}
If the data are $\Normal(\mu,\sigma^2)$, then $(\Xbar-\mu)/(\sigma/\sqrt n)$ is
exactly $\Normal(0,1)$, $(n-1)s^2/\sigma^2$ has a chi-square distribution with $n-1$
degrees of freedom, and the two are independent. A standard Normal divided by the
square root of an independent (chi-square / its degrees of freedom) is, by
definition, $t_{n-1}$. The $\sigma$ cancels, which is why the test needs no knowledge
of it. (Derivation: see Probability From the Ground Up, Ch.~32.)
\end{whybox}

\begin{example}[Does the service meet its claim?]
A vendor claims mean response time $120$ ms. You time $n=10$ requests:
\[
  128,\ 135,\ 119,\ 142,\ 126,\ 131,\ 124,\ 138,\ 117,\ 133 .
\]
\textbf{Summaries.} The sum is $1293$, so $\xbar = 129.3$. The squared deviations from
$129.3$ sum to $584.1$, so $s^2 = 584.1/9 = 64.9$ and $s = 8.056$. The standard
error is $s/\sqrt{10} = 2.548$.

\textbf{Statistic.} $H_0\colon\mu=120$ versus $H_1\colon \mu\ne120$:
\[
  t = \frac{129.3 - 120}{2.548} = 3.65, \qquad \text{df} = 9 .
\]
\textbf{p-value.} $p = \P(|T_9| \ge 3.65) = 0.0053$. The cut-off for $\alpha=0.05$ is
$t_{9,\,0.975}=2.262$, and $3.65$ is well past it. Reject: the true mean is very
likely above $120$ ms.
\end{example}

\begin{figure}[H]
\centering
\begin{tikzpicture}[x=0.42cm]
\draw[->] (0,0) -- (26.5,0) node[right, font=\small] {ms};
\foreach \t in {115,120,...,140} {
  \draw (\t-115,0.1) -- (\t-115,-0.1) node[below, font=\small] {\t};
}
\foreach \v in {128,135,119,142,126,131,124,138,117,133}
  \fill[distblue] (\v-115,0.35) circle (2.6pt);
\draw[trapred, thick, dashed] (5,-0.1) -- (5,1.5) node[above, font=\small] {$\mu_0 = 120$};
\draw[targetgreen!60!black, very thick] (14.3,-0.1) -- (14.3,1.5) node[above, font=\small] {$\bar x = 129.3$};
\draw[{Stealth}-{Stealth}, structurepurple] (5,0.95) -- (14.3,0.95);
\node[font=\footnotesize, structurepurple, fill=white, inner sep=1pt] at (9.65,0.95) {gap $9.3$ ms $=3.65\,\SE$};
\end{tikzpicture}
\caption{The ten response times (one dot per request), the claimed mean and the
sample mean. The gap of
$9.3$ ms is $3.65$ standard errors: large relative to how much $\bar x$ wobbles from
sample to sample, even though it is small relative to the spread of individual
requests.}
\label{fig:ht-onesample}
\end{figure}

\section{Two samples: pooled and Welch t-tests}
\label{sec:ht-t2}

Return to the render times. Now test $H_0\colon \mu_{\text{old}} = \mu_{\text{new}}$
using a formula instead of shuffles. The statistic is again
``difference over its standard error'':
\[
  t = \frac{\bar x_1 - \bar x_2}{\SE(\bar x_1 - \bar x_2)} .
\]
The two versions of the test differ only in how they estimate that standard error.

\begin{figure}[H]
\centering
\begin{tikzpicture}[x=0.15cm]
\draw[->] (0,0) -- (72,0) node[right, font=\small] {ms};
\foreach \t in {175,185,...,245} \draw (\t-175,0.1) -- (\t-175,-0.1) node[below, font=\small] {\t};
\foreach \v in {212,198,225,240,205,231} \fill[sourceblue] (\v-175,2.0) circle (2.8pt);
\foreach \v in {190,184,207,195,178,201} \fill[targetgreen] (\v-175,0.6) circle (2.8pt);
\draw[sourceblue, very thick] (43.5,1.75) -- (43.5,2.25);
\draw[targetgreen, very thick] (17.5,0.35) -- (17.5,0.85);
\node[anchor=east, font=\small, sourceblue] at (-0.5,2.0) {old};
\node[anchor=east, font=\small, targetgreen!70!black] at (-0.5,0.6) {new};
\node[anchor=south, font=\footnotesize, sourceblue] at (43.5,2.3) {$\bar x_1=218.5,\ s_1=16.2$};
\node[anchor=south, font=\footnotesize, targetgreen!70!black] at (17.5,0.95) {$\bar x_2=192.5,\ s_2=10.7$};
\draw[{Stealth}-{Stealth}, structurepurple, thick] (17.5,-1.05) -- (43.5,-1.05);
\draw[structurepurple, densely dotted] (17.5,0.35) -- (17.5,-1.2);
\draw[structurepurple, densely dotted] (43.5,1.75) -- (43.5,-1.2);
\node[structurepurple, font=\footnotesize, anchor=north] at (30.5,-1.15) {difference in means $26.0$ ms};
\end{tikzpicture}
\caption{Render times by version (one dot per page load), with group means as
vertical ticks. The groups overlap a little, and the old group is more spread out.
The t-test asks whether the $26$ ms gap between the means (purple arrow) is large
compared with this within-group noise.}
\label{fig:ht-twosample}
\end{figure}

\textbf{Pooled (Student) t-test.} Assume both groups share one variance $\sigma^2$.
Estimate it by combining the two sample variances, weighted by their degrees of
freedom:
\[
  s_p^2 = \frac{(n_1-1)s_1^2 + (n_2-1)s_2^2}{n_1+n_2-2}, \qquad
  \SE = s_p\sqrt{\frac{1}{n_1}+\frac{1}{n_2}},
  \qquad \text{df} = n_1+n_2-2 .
\]
For the render times, $s_1^2 = 261.1$ and $s_2^2 = 115.5$:
\[
  s_p^2 = \frac{5(261.1)+5(115.5)}{10} = 188.3, \quad
  \SE = 13.72\sqrt{\tfrac16+\tfrac16} = 7.92, \quad
  t = \frac{26.0}{7.92} = 3.28 .
\]
With $10$ degrees of freedom, $p = 0.0083$.

\textbf{Welch t-test.} Drop the equal-variance assumption. Each mean has its own
variance $s_i^2/n_i$, and variances of independent quantities add:
\[
  \SE = \sqrt{\frac{s_1^2}{n_1}+\frac{s_2^2}{n_2}}, \qquad
  \text{df} \approx \frac{\left(s_1^2/n_1 + s_2^2/n_2\right)^2}
  {\dfrac{(s_1^2/n_1)^2}{n_1-1} + \dfrac{(s_2^2/n_2)^2}{n_2-1}} .
\]
The odd-looking df formula (the Welch--Satterthwaite approximation) picks the $t$
distribution whose spread best matches the actual null distribution. It always lies
between $\min(n_1,n_2)-1$ and $n_1+n_2-2$. Here $n_1=n_2$, so the SE is again
$\sqrt{261.1/6+115.5/6} = 7.92$ and $t = 3.28$, but df $= 8.70$, so
$p = 0.0099$.

\begin{center}
\begin{tabular}{lccc}
\toprule
Method & assumes & statistic & two-sided $p$\\
\midrule
Permutation (exact) & exchangeable labels & $26.0$ ms & $0.0108$\\
Pooled t-test & Normal, equal variances & $t=3.28$, df $=10$ & $0.0083$\\
Welch t-test & Normal & $t=3.28$, df $=8.70$ & $0.0099$\\
\bottomrule
\end{tabular}
\end{center}

Three different routes to the null world, three nearly identical answers. The
permutation test is a useful sanity check on the formulas, and the formulas explain
why the permutation histogram in Figure~\ref{fig:ht-perm-null} looks like a bell.

\begin{trapbox}
Use Welch by default. When variances really are equal it loses almost nothing. When
they are not, and the group sizes differ, the pooled test can be badly wrong
(Section~\ref{sec:ht-robust} shows a case where its false-alarm rate is $30\%$
instead of $5\%$). Pre-testing for equal variances and then choosing is worse than
just using Welch.
\end{trapbox}

\section{Paired data}
\label{sec:ht-paired}

A build engineer measures a test suite's runtime (seconds) on $8$ different CI
machines, once with the old compiler and once with the new one:
\begin{center}
\small
\begin{tabular}{lcccccccc}
\toprule
machine & 1 & 2 & 3 & 4 & 5 & 6 & 7 & 8\\
\midrule
old $a_i$ & 52.1 & 61.4 & 47.8 & 70.2 & 55.6 & 66.3 & 49.9 & 58.7\\
new $b_i$ & 50.8 & 59.9 & 47.1 & 68.1 & 54.2 & 65.0 & 48.3 & 57.9\\
$d_i=a_i-b_i$ & 1.3 & 1.5 & 0.7 & 2.1 & 1.4 & 1.3 & 1.6 & 0.8\\
\bottomrule
\end{tabular}
\end{center}

Pretend for a moment these were $16$ unrelated runs and apply Welch. The means are
$57.75$ and $56.41$, a gap of $1.34$ s. But the machines differ wildly (from $47$ s
to $70$ s), so $s_a = 7.92$, $s_b = 7.69$, $\SE = 3.90$, $t = 0.34$, and $p=0.74$.
No evidence of anything.

Now look at the differences. Every single machine got faster, by between $0.7$ and
$2.1$ seconds. Treat the $8$ differences as one sample and run a one-sample t-test of
$H_0\colon \mu_d = 0$:
\[
  \bar d = 1.3375, \quad s_d = 0.4438, \quad
  \SE = \frac{0.4438}{\sqrt 8} = 0.1569, \quad
  t = \frac{1.3375}{0.1569} = 8.52, \quad \text{df}=7 .
\]
$p = 0.00006$. Same data, opposite conclusion. The paired analysis is the correct
one, because the design paired each old run with a new run on the same machine.

\begin{figure}[H]
\centering
\begin{tikzpicture}
\begin{axis}[
  name=slope, width=5.6cm, height=6cm, xmin=-0.3, xmax=1.3, ymin=45, ymax=72,
  xtick={0,1}, xticklabels={old,new}, ylabel={runtime (s)}, axis lines=left,
  title={Raw runtimes, paired}, title style={font=\small},
]
\addplot[mark=*, mark size=1.6pt, distblue] coordinates {(0,52.1)(1,50.8)};
\addplot[mark=*, mark size=1.6pt, distblue] coordinates {(0,61.4)(1,59.9)};
\addplot[mark=*, mark size=1.6pt, distblue] coordinates {(0,47.8)(1,47.1)};
\addplot[mark=*, mark size=1.6pt, distblue] coordinates {(0,70.2)(1,68.1)};
\addplot[mark=*, mark size=1.6pt, distblue] coordinates {(0,55.6)(1,54.2)};
\addplot[mark=*, mark size=1.6pt, distblue] coordinates {(0,66.3)(1,65.0)};
\addplot[mark=*, mark size=1.6pt, distblue] coordinates {(0,49.9)(1,48.3)};
\addplot[mark=*, mark size=1.6pt, distblue] coordinates {(0,58.7)(1,57.9)};
\end{axis}
\begin{axis}[
  at={(slope.east)}, anchor=west, xshift=1.2cm,
  width=6.3cm, height=6cm, xmin=-0.3, xmax=2.5, ymin=0, ymax=3.5,
  ytick=\empty, axis x line=bottom, axis y line=none, xlabel={$d_i = a_i - b_i$ (s)},
  title={Differences}, title style={font=\small}, clip=false,
]
\addplot[only marks, mark=*, mark size=2.4pt, quotientorange] coordinates {
(0.7,0.4)(0.8,0.4)(1.3,0.4)(1.3,0.8)(1.4,0.4)(1.5,0.4)(1.6,0.4)(2.1,0.4)};
\draw[trapred, dashed, thick] (axis cs:0,0) -- (axis cs:0,2.6) node[above, font=\footnotesize] {$H_0$: $0$};
\draw[targetgreen!60!black, very thick] (axis cs:1.3375,1.0) -- (axis cs:1.3375,1.6);
\node[targetgreen!60!black, font=\footnotesize, anchor=south] at (axis cs:1.3375,1.6) {$\bar d=1.34$};
\draw[{Stealth}-{Stealth}, targetgreen!60!black] (axis cs:1.0237,1.25) -- (axis cs:1.6513,1.25);
\node[font=\footnotesize, anchor=west, targetgreen!60!black] at (axis cs:1.67,1.25) {$\pm 2\,\SE$};
\end{axis}
\end{tikzpicture}
\caption{Left: each line is one machine. Machine-to-machine differences (the vertical
spread) are huge, but every line slopes down. Right: the differences remove the
machine effect and cluster tightly around $1.34$ s, far from $0$.}
\label{fig:ht-paired}
\end{figure}

\begin{whybox}
For any two random quantities,
\[
  \Var(A - B) = \Var(A) + \Var(B) - 2\Cov(A,B).
\]
An unpaired analysis implicitly sets the covariance to zero. But a slow machine is
slow under both compilers, so $A$ and $B$ are strongly positively correlated (here
$r = 0.9988$). The covariance term then cancels almost all the variance:
$7.92^2 + 7.69^2 - 2(0.9988)(7.92)(7.69) \approx 0.197 = 0.4438^2$. Pairing subtracts
the machine effect before it can drown the compiler effect. This is why
before/after studies, matched pairs and ``same user, both variants'' designs are so
much more sensitive than comparing two separate groups.
\end{whybox}

\begin{recoverybox}
Pause and think: an A/B test assigns each user to only one variant. Paired or two
sample? (Two sample: no user appears in both groups, so there is nothing to pair.)
\end{recoverybox}

\section{Errors, significance and power}
\label{sec:ht-power}

A test ends in a decision, and a decision can be wrong in two ways.

\begin{center}
\begin{tabular}{lcc}
\toprule
 & $H_0$ true (nothing going on) & $H_0$ false (real effect)\\
\midrule
Reject $H_0$ & \textbf{Type I error} (false alarm), prob.\ $\alpha$ & correct, prob.\ $1-\beta$ = \textbf{power}\\
Do not reject & correct, prob.\ $1-\alpha$ & \textbf{Type II error} (miss), prob.\ $\beta$\\
\bottomrule
\end{tabular}
\end{center}

$\alpha$ is the false-alarm rate you chose. $\beta$ (beta) is the probability of
missing a real effect, and \textbf{power} $=1-\beta$ is the probability of detecting
it. Unlike $\alpha$, power is not a single number: it depends on how big the real
effect is.

To compute power you need the picture the null hypothesis alone cannot give: an
\emph{alternative world}, as in panel 6 of Figure~\ref{fig:ht-flow}. Suppose the
coin friend truly is right $75\%$ of the time.
Her score is then $\Bin(20, 0.75)$. Overlay that world on the null world and keep the
same rejection rule, $X \ge 15$ (Figure~\ref{fig:ht-coin-power}).

\begin{figure}[H]
\centering
\begin{tikzpicture}
\begin{axis}[
  width=12cm, height=5.6cm, xlabel={correct calls out of 20}, ylabel={probability},
  xmin=1.5, xmax=20.5, ymin=0, ymax=0.23, xtick={2,4,...,20}, axis lines=left,
  legend style={font=\footnotesize, draw=none, at={(0.5,-0.3)}, anchor=north}, legend columns=2,
  yticklabel style={/pgf/number format/fixed}, clip=false,
]
\addplot[ybar, bar width=0.3cm, bar shift=-0.16cm, area legend, fill=histfill, draw=distblue] coordinates {
(2,0.0002)(3,0.0011)(4,0.0046)(5,0.0148)(6,0.0370)(7,0.0739)(8,0.1201)(9,0.1602)
(10,0.1762)(11,0.1602)(12,0.1201)(13,0.0739)(14,0.0370)};
\addlegendentry{null $\theta=0.5$: do not reject}
\addplot[ybar, bar width=0.3cm, bar shift=-0.16cm, area legend, fill=trapred, draw=trapred!70!black] coordinates {
(15,0.0148)(16,0.0046)(17,0.0011)(18,0.0002)};
\addlegendentry{null: false alarm, $\alpha=0.021$}
\addplot[ybar, bar width=0.3cm, bar shift=0.16cm, area legend, fill=bootorange!50, draw=bootorange] coordinates {
(7,0.0002)(8,0.0008)(9,0.0030)(10,0.0099)(11,0.0271)(12,0.0609)(13,0.1124)(14,0.1686)};
\addlegendentry{alt.\ $\theta=0.75$: miss, $\beta=0.383$}
\addplot[ybar, bar width=0.3cm, bar shift=0.16cm, area legend, fill=targetgreen!60, draw=targetgreen!60!black] coordinates {
(15,0.2023)(16,0.1897)(17,0.1339)(18,0.0669)(19,0.0211)(20,0.0032)};
\addlegendentry{alt.: detected, power $=0.617$}
\draw[densityred, dashed, thick] (axis cs:14.5,0) -- (axis cs:14.5,0.22);
\end{axis}
\end{tikzpicture}
\caption{Null world (left bars of each pair) and an alternative world where her true
hit rate is $0.75$ (right bars). The dashed line is the rejection threshold. Red is
the false-alarm probability; green is the power. Even a genuinely skilled friend
fails this $20$-flip test $38\%$ of the time.}
\label{fig:ht-coin-power}
\end{figure}

\subsection{The general picture}

For a z-type statistic the two worlds are two bells. Under $H_0$ the statistic is
$\Normal(0,1)$. Under an alternative with true effect $\delta$ it is $\Normal(\delta/\SE,
1)$: the same bell, shifted by the effect measured in standard errors.
Figure~\ref{fig:ht-power} shows a shift of $2.8$.

\begin{figure}[H]
\centering
\begin{tikzpicture}
\begin{axis}[
  width=12cm, height=5.6cm, domain=-3.5:6.5, samples=160, axis lines=left,
  ymin=0, ymax=0.5, xlabel={test statistic $z$}, ytick=\empty,
  xtick={0,1.96,2.8}, xticklabels={$0$,$1.96$,$2.8$}, clip=false,
]
\addplot[fill=bootorange!45, draw=none, domain=-1:1.96] {0.39894*exp(-0.5*(x-2.8)^2)} \closedcycle;
\addplot[fill=targetgreen!35, draw=none, domain=1.96:6.5] {0.39894*exp(-0.5*(x-2.8)^2)} \closedcycle;
\addplot[fill=trapred!70, draw=none, domain=1.96:3.5] {0.39894*exp(-0.5*x^2)} \closedcycle;
\addplot[fill=trapred!70, draw=none, domain=-3.5:-1.96] {0.39894*exp(-0.5*x^2)} \closedcycle;
\addplot[very thick, distblue] {0.39894*exp(-0.5*x^2)};
\addplot[very thick, targetgreen!60!black] {0.39894*exp(-0.5*(x-2.8)^2)};
\draw[densityred, dashed, thick] (axis cs:1.96,0) -- (axis cs:1.96,0.46);
\node[distblue, font=\small, anchor=south] at (axis cs:-0.6,0.41) {$H_0$ world};
\node[targetgreen!60!black, font=\small, anchor=south] at (axis cs:3.6,0.41) {$H_1$ world};
\node[font=\small, anchor=west, targetgreen!50!black] at (axis cs:4.5,0.22) {power $=0.80$};
\node[font=\small, anchor=south, bootorange!80!black] at (axis cs:1.0,0.45) {$\beta=0.20$};
\draw[bootorange!80!black] (axis cs:1.0,0.45) -- (axis cs:1.6,0.05);
\node[font=\small, anchor=south, trapred, align=center] at (axis cs:-2.5,0.2) {$\alpha/2=0.025$\\each tail};
\draw[trapred] (axis cs:-2.5,0.2) -- (axis cs:-2.2,0.012);
\end{axis}
\end{tikzpicture}
\caption{Two overlapping sampling distributions. Blue: the statistic if $H_0$ is true.
Green: the statistic if the true effect is $2.8$ standard errors. Red tails: $\alpha$.
Orange: $\beta$, the part of the $H_1$ world that falls short of the cut-off. Green
fill: power. Moving the cut-off right lowers $\alpha$ but raises $\beta$.}
\label{fig:ht-power}
\end{figure}

Reading the picture gives the power formula. Ignoring the tiny left tail,
\[
  \text{power} \approx \P\!\left(Z > z_{\alpha/2} - \frac{\delta}{\SE}\right)
  = \Phi\!\left(\frac{\delta}{\SE} - z_{\alpha/2}\right),
\]
where $z_{\alpha/2}$ is the cut-off ($1.96$ for $\alpha = 0.05$). With $\delta/\SE =
2.8$: $\Phi(2.8-1.96) = \Phi(0.84) = 0.80$. Three levers raise power:
\begin{itemize}
\item a larger effect $\delta$ (you rarely control this);
\item a smaller $\SE$, which means more data ($\SE\propto 1/\sqrt n$), less noise, or
a better design such as pairing;
\item a larger $\alpha$ (more false alarms, the trade-off in the figure).
\end{itemize}

For comparing two groups of $n$ each with common SD $\sigma$, $\SE = \sigma\sqrt{2/n}$.
Write the effect as $d = \delta/\sigma$, the difference in units of SD (``Cohen's
$d$'', a scale-free effect size: $0.2$ is small, $0.5$ medium, $0.8$ large; see
Chapter 11). Then $\delta/\SE = d\sqrt{n/2}$ and
\[
  \text{power} \approx \Phi\!\left(d\sqrt{n/2} - 1.96\right).
\]

\begin{figure}[H]
\centering
\begin{tikzpicture}
\begin{axis}[
  width=9.6cm, height=6cm, xmin=0, xmax=200, ymin=0, ymax=1.05,
  xlabel={sample size per group $n$}, ylabel={power}, axis lines=left,
  grid=major, grid style={gray!20}, clip=false,
]
\addplot[thick, trapred] coordinates {(2,0.055) (5,0.062) (10,0.073) (15,0.085) (20,0.097) (25,0.109) (30,0.121) (40,0.145) (50,0.170) (60,0.195) (70,0.219) (80,0.244) (100,0.293) (120,0.341) (140,0.387) (160,0.432) (180,0.475) (200,0.516)};
\addplot[thick, quotientorange] coordinates {(2,0.079) (5,0.124) (10,0.201) (15,0.278) (20,0.353) (25,0.424) (30,0.491) (40,0.609) (50,0.705) (60,0.782) (70,0.841) (80,0.885) (100,0.942) (120,0.972) (140,0.987) (160,0.994) (180,0.997) (200,0.999)};
\addplot[thick, distblue] coordinates {(2,0.126) (5,0.244) (10,0.432) (15,0.591) (20,0.716) (25,0.807) (30,0.873) (40,0.947) (50,0.979) (60,0.992) (70,0.997) (80,0.999) (100,1.000) (120,1.000) (140,1.000) (160,1.000) (180,1.000) (200,1.000)};
\node[distblue, font=\small, anchor=west] at (axis cs:202,1.04) {large, $d=0.8$};
\node[quotientorange!80!black, font=\small, anchor=west] at (axis cs:202,0.93) {medium, $d=0.5$};
\node[trapred, font=\small, anchor=west] at (axis cs:202,0.516) {small, $d=0.2$};
\draw[gray, dashed] (axis cs:0,0.8) -- (axis cs:200,0.8);
\node[gray!70!black, font=\footnotesize, anchor=north east] at (axis cs:200,0.79) {target power $0.8$};
\draw[gray, dashed] (axis cs:0,0.05) -- (axis cs:200,0.05);
\node[gray!70!black, font=\footnotesize, anchor=south east] at (axis cs:200,0.05) {$\alpha=0.05$};
\fill[distblue] (axis cs:25,0.8) circle (2pt);
\node[distblue, font=\footnotesize, anchor=north west] at (axis cs:27,0.78) {$n=25$};
\fill[quotientorange] (axis cs:63,0.8) circle (2pt);
\node[quotientorange!80!black, font=\footnotesize, anchor=north west] at (axis cs:64,0.78) {$n=63$};
\node[trapred, font=\footnotesize, anchor=north west, align=left] at (axis cs:95,0.66) {small effect reaches $0.8$\\only at $n\approx 390$ (off chart)};
\end{axis}
\end{tikzpicture}
\caption{Power of a two-sided, two-sample test at $\alpha=0.05$ as the sample size
grows (Normal approximation). The dots mark where each curve reaches $80\%$ power:
about $25$ per group for a large effect, $63$ for a medium one, and about $390$ for a
small one. As $n$ shrinks towards $0$, power falls towards $\alpha$: you would
reject only by chance.}
\label{fig:ht-power-n}
\end{figure}

\begin{figure}[H]
\centering
\begin{tikzpicture}
\begin{axis}[
  width=11.5cm, height=6cm, xmin=0, xmax=1.2, ymin=0, ymax=1.05,
  xlabel={true effect size $d$}, ylabel={power}, axis lines=left,
  legend style={font=\small}, legend pos=south east, legend cell align=left, grid=major, grid style={gray!20},
  clip=false,
]
\addplot[thick, trapred] coordinates {(0.00,0.050) (0.05,0.053) (0.10,0.062) (0.15,0.076) (0.20,0.097) (0.25,0.124) (0.30,0.158) (0.35,0.198) (0.40,0.244) (0.45,0.296) (0.50,0.353) (0.55,0.413) (0.60,0.475) (0.65,0.538) (0.70,0.600) (0.75,0.660) (0.80,0.716) (0.85,0.767) (0.90,0.812) (0.95,0.852) (1.00,0.885) (1.05,0.913) (1.10,0.936) (1.15,0.953) (1.20,0.967)};
\addlegendentry{$n=20$ per group}
\addplot[thick, quotientorange] coordinates {(0.00,0.050) (0.05,0.057) (0.10,0.079) (0.15,0.117) (0.20,0.170) (0.25,0.240) (0.30,0.323) (0.35,0.417) (0.40,0.516) (0.45,0.614) (0.50,0.705) (0.55,0.785) (0.60,0.851) (0.65,0.901) (0.70,0.938) (0.75,0.963) (0.80,0.979) (0.85,0.989) (0.90,0.994) (0.95,0.997) (1.00,0.999) (1.05,0.999) (1.10,1.000) (1.15,1.000) (1.20,1.000)};
\addlegendentry{$n=50$}
\addplot[thick, distblue] coordinates {(0.00,0.050) (0.05,0.079) (0.10,0.170) (0.15,0.323) (0.20,0.516) (0.25,0.705) (0.30,0.851) (0.35,0.938) (0.40,0.979) (0.45,0.994) (0.50,0.999) (0.55,1.000) (0.60,1.000) (0.65,1.000) (0.70,1.000) (0.75,1.000) (0.80,1.000) (0.85,1.000) (0.90,1.000) (0.95,1.000) (1.00,1.000) (1.05,1.000) (1.10,1.000) (1.15,1.000) (1.20,1.000)};
\addlegendentry{$n=200$}
\draw[gray!70!black, dashed] (axis cs:0.3,0) -- (axis cs:0.3,1.0);
\fill[distblue] (axis cs:0.3,0.851) circle (2pt);
\node[distblue, font=\footnotesize, anchor=east] at (axis cs:0.29,0.87) {$0.85$};
\fill[trapred] (axis cs:0.3,0.158) circle (2pt);
\node[trapred, font=\footnotesize, anchor=west] at (axis cs:0.31,0.13) {$0.16$};
\node[gray!70!black, font=\footnotesize, anchor=south] at (axis cs:0.3,1.0) {true effect $d=0.3$};
\end{axis}
\end{tikzpicture}
\caption{Power curves: the probability of rejecting $H_0$ as a function of the true
effect. At $d=0$ (the null is true) every curve equals $\alpha=0.05$. Bigger samples
give steeper curves. Follow the dashed line at $d=0.3$: $n=200$ per group detects it
with $85\%$ power, while $n=20$ detects it only $16\%$ of the time.}
\label{fig:ht-power-d}
\end{figure}

\subsection{Sample size for an A/B test}

Power calculations are done \emph{before} the experiment, to decide how much data to
collect. Here is the standard calculation.

\begin{example}[How many users?]
A checkout page converts $10\%$ of visitors. A redesign might lift that to $11\%$.
How many users per variant do you need to detect that lift with $\alpha=0.05$
(two-sided) and $80\%$ power?

\textbf{Step 1: the SE.} With $n$ users per arm, the observed conversion rates
$\hat p_1,\hat p_2$ have variances $p(1-p)/n$ (Chapter 6), so
\[
  \SE(\hat p_2-\hat p_1) = \sqrt{\frac{p_1(1-p_1) + p_2(1-p_2)}{n}} .
\]
\textbf{Step 2: the requirement.} From Figure~\ref{fig:ht-power}, power $1-\beta$
requires the alternative bell to sit $z_{\alpha/2} + z_\beta$ standard errors from
zero, where $z_\beta$ is the point with area $\beta$ above it ($0.8416$ for
$\beta=0.2$). So $\delta/\SE = z_{\alpha/2}+z_\beta$. Solving for $n$:
\[
  n = \frac{(z_{\alpha/2} + z_\beta)^2\,\bigl[p_1(1-p_1)+p_2(1-p_2)\bigr]}{\delta^2}.
\]
\textbf{Step 3: the numbers.} $z_{\alpha/2}+z_\beta = 1.960+0.842 = 2.802$, whose
square is $7.849$. Next, $p_1(1-p_1)+p_2(1-p_2) = 0.0900 + 0.0979 = 0.1879$, and
$\delta^2 = 0.01^2 = 0.0001$. So
\[
  n = \frac{7.849 \times 0.1879}{0.0001} \approx 14{,}748 \text{ users per arm},
\]
about $29{,}500$ in total.
\end{example}

Two features of the formula drive every A/B planning conversation. First,
$n\propto 1/\delta^2$: detecting a lift half as big ($0.5$ points) needs about four
times the users, $57{,}761$ per arm (a little under four times $14{,}748$, because
$p_2(1-p_2)$ also shrinks slightly). Second, the variance term $p(1-p)$ is largest
near $50\%$, so low-conversion metrics are cheaper per user but their absolute lifts
are tiny. Figure~\ref{fig:ht-ab-n} shows how steep the first effect is.

\begin{figure}[H]
\centering
\begin{tikzpicture}
\begin{axis}[
  width=11cm, height=6cm, ymode=log, xmin=0, xmax=3.2, ymin=1000, ymax=400000,
  domain=0.25:3, samples=80, xtick={0.25,0.5,1,1.5,2,2.5,3},
  xticklabels={0.25,0.5,1,1.5,2,2.5,3},
  ytick={1000,10000,100000}, yticklabels={1{,}000,10{,}000,100{,}000},
  xlabel={lift to detect (percentage points above a 10\% baseline)},
  ylabel={users needed per arm (log scale)}, axis lines=left,
  grid=major, grid style={gray!20}, clip=false,
]
\addplot[very thick, distblue] {7.849*(0.09+(0.1+x/100)*(0.9-x/100))/(x/100)^2};
\fill[trapred] (axis cs:1,14748) circle (2.5pt);
\node[trapred, font=\small, anchor=south west] at (axis cs:1.05,15500) {1 point: $14{,}748$};
\fill[trapred] (axis cs:0.5,57761) circle (2.5pt);
\node[trapred, font=\small, anchor=south west] at (axis cs:0.55,60000) {0.5 point: $57{,}761$};
\fill[trapred] (axis cs:2,3838) circle (2.5pt);
\node[trapred, font=\small, anchor=south west] at (axis cs:2.05,4000) {2 points: $3{,}838$};
\end{axis}
\end{tikzpicture}
\caption{Users per arm for $\alpha=0.05$ (two-sided) and $80\%$ power, baseline
conversion $10\%$. Halving the lift you want to detect roughly quadruples the users
needed; the vertical axis is logarithmic, so each grid line is ten times the one below.}
\label{fig:ht-ab-n}
\end{figure}

\begin{lstlisting}[language=Python]
from scipy.stats import norm
p1, p2, alpha, power = 0.10, 0.11, 0.05, 0.80
z = norm.ppf(1 - alpha/2) + norm.ppf(power)            # 1.960 + 0.842
n = z**2 * (p1*(1-p1) + p2*(1-p2)) / (p2 - p1)**2
print(round(n))                                        # 14748 per arm
\end{lstlisting}

\begin{trapbox}
An underpowered study is not a cheap study; it is a wasted one. With $20\%$ power,
four out of five real effects come back ``not significant'', and the ones that do
reach significance tend to overstate the effect (only lucky, inflated estimates clear
the bar). Always run the power calculation first.
\end{trapbox}

\begin{historybox}
In the 1920s R.\,A.\ Fisher, working on crop experiments at Rothamsted in England,
needed a way to say whether a fertilizer's apparent benefit could be due to chance
variation between field plots. He introduced the p-value as a graded measure of
surprise under a null hypothesis, and popularized $0.05$ as a convenient bar. His
famous ``lady tasting tea'' example (1935), in which a woman claims she can tell
whether milk was poured before the tea, is the first permutation test: count the
ways the cups could be labelled by pure guessing. In 1933 Jerzy Neyman and Egon
Pearson reframed testing as a decision rule. They added the alternative hypothesis,
the two error types, and power, which made it possible to design an experiment
with a known chance of finding an effect. Modern practice is a blend of the two
views, which explains some of its confusions.
\end{historybox}

\section{Tests and confidence intervals are two views of one thing}
\label{sec:ht-duality}

Take the response-time data again: $\xbar = 129.3$, $\SE = 2.548$, df $=9$. The $95\%$
confidence interval from Chapter 7 is
\[
  129.3 \pm 2.262 \times 2.548 = 129.3 \pm 5.76 = [123.54,\ 135.06].
\]
Now test $H_0\colon\mu=\mu_0$ for several values of $\mu_0$ at $\alpha=0.05$:

\begin{center}
\begin{tabular}{cccc}
\toprule
$\mu_0$ & $t$ & $p$ & decision\\
\midrule
$120$ & $3.65$ & $0.005$ & reject\\
$125$ & $1.69$ & $0.126$ & do not reject\\
$130$ & $-0.27$ & $0.790$ & do not reject\\
$137$ & $-3.02$ & $0.014$ & reject\\
\bottomrule
\end{tabular}
\end{center}

\begin{figure}[H]
\centering
\begin{tikzpicture}[x=0.42cm]
\draw[->] (0,0) -- (26.5,0) node[right, font=\small] {$\mu$ (ms)};
\foreach \t in {115,120,...,140} \draw (\t-115,0.1) -- (\t-115,-0.1) node[below, font=\small] {\t};
\draw[targetgreen!60!black, line width=3pt] (8.54,0.8) -- (20.06,0.8);
\fill[targetgreen!60!black] (14.3,0.8) circle (3.5pt);
\node[font=\small, targetgreen!50!black, anchor=south] at (14.3,1.0) {95\% CI $[123.5,\,135.1]$};
\foreach \m/\c/\lab in {5/trapred/{$\mu_0=120$\\rejected}, 10/sourceblue/{$\mu_0=125$\\kept}, 15/sourceblue/{$\mu_0=130$\\kept}, 22/trapred/{$\mu_0=137$\\rejected}} {
  \draw[\c, thick, -{Stealth}] (\m,-1.6) -- (\m,-0.75);
  \node[\c, font=\footnotesize, anchor=north, align=center] at (\m,-1.65) {\lab};
}
\end{tikzpicture}
\caption{Duality between tests and intervals. The green bar is the $95\%$ CI; the
arrows are four candidate null values. The two inside it ($125$ and $130$) are kept
at $\alpha=0.05$; the two outside it ($120$ and $137$) are rejected. The CI is exactly
the set of null values the test would not reject.}
\label{fig:ht-duality}
\end{figure}

\begin{theorem}[Duality]
A two-sided test at level $\alpha$ rejects $H_0\colon\mu=\mu_0$ exactly when $\mu_0$ lies
outside the $100(1-\alpha)\%$ confidence interval built from the same statistic.
\end{theorem}

\begin{whybox}
The test does not reject when $|\xbar-\mu_0|/\SE \le t^*$, where $t^*$ is the
cut-off. Rearranged, that is $\xbar - t^*\SE \le \mu_0 \le \xbar + t^*\SE$, which is
the confidence interval. Same inequality, solved for a different variable.
\end{whybox}

This is why the interval is usually the better report. A p-value answers one
question (``is $\mu=120$ plausible?''); the interval answers it for every $\mu_0$ at
once, and also shows how big the effect is and how precisely it is known.

\section{What a p-value is not}
\label{sec:ht-not}

\begin{trapbox}
A p-value is $\P(\text{data this extreme} \mid H_0)$. It is \emph{not}:
\begin{enumerate}
\item \textbf{The probability that $H_0$ is true.} $p=0.03$ does not mean ``$3\%$
chance of no effect''. That would be $\P(H_0\mid\text{data})$, a different conditional
probability that needs a prior (Chapter 14).
\item \textbf{The probability the result is a fluke.} Same mistake in different words.
\item \textbf{One minus the probability that $H_1$ is true.} $p=0.01$ does not make
the alternative $99\%$ likely.
\item \textbf{A measure of effect size.} A trivial $0.1$ ms speedup gives a tiny p-value
with enough data (Figure~\ref{fig:ht-absence}, right). Report the effect and its CI.
\item \textbf{Evidence for $H_0$ when large.} $p=0.6$ means ``compatible with $H_0$'',
not ``$H_0$ shown'' (Section~\ref{sec:ht-absence}).
\item \textbf{A bright line.} $p=0.049$ and $p=0.051$ are essentially the same
evidence.
\item \textbf{The probability of replicating.} A study with $p=0.05$ has only about a
$50\%$ chance that an exact repeat also gives $p<0.05$, if the true effect equals the
observed one.
\end{enumerate}
\end{trapbox}

Item 7 surprises most people, and a picture shows why. Suppose the original study
landed exactly on the line, $z=1.96$ ($p=0.05$), and the true effect really is that
size. An exact repeat then produces a $z$ that wobbles around $1.96$ with standard
deviation $1$ (Figure~\ref{fig:ht-replicate}). Half of that bell lies below the cut-off.

\begin{figure}[H]
\centering
\begin{tikzpicture}
\begin{axis}[
  width=11.5cm, height=5cm, domain=-1.5:5.5, samples=140, axis lines=left,
  ymin=0, ymax=0.5, ytick=\empty, xmin=-1.5, xmax=5.5,
  xtick={0,1.96,4}, xticklabels={$0$,$1.96$,$4$},
  xlabel={$z$ in the exact repeat study}, ylabel={density}, clip=false,
]
\addplot[fill=bootorange!45, draw=none, domain=-1.5:1.96] {0.39894*exp(-0.5*(x-1.96)^2)} \closedcycle;
\addplot[fill=targetgreen!40, draw=none, domain=1.96:5.5] {0.39894*exp(-0.5*(x-1.96)^2)} \closedcycle;
\addplot[very thick, targetgreen!60!black] {0.39894*exp(-0.5*(x-1.96)^2)};
\draw[densityred, dashed, thick] (axis cs:1.96,0) -- (axis cs:1.96,0.46);
\node[densityred, font=\small, anchor=south] at (axis cs:1.96,0.46) {cut-off $1.96$ $=$ original study's $z$};
\node[bootorange!80!black, font=\small, align=center] (miss) at (axis cs:-0.5,0.3) {repeat misses\\$p\ge 0.05$: $50\%$};
\node[targetgreen!50!black, font=\small, align=center] (hit) at (axis cs:4.5,0.3) {repeat succeeds\\$p<0.05$: $50\%$};
\draw[bootorange!80!black] (miss.south) -- (axis cs:1.2,0.12);
\draw[targetgreen!50!black] (hit.south) -- (axis cs:2.7,0.12);
\end{axis}
\end{tikzpicture}
\caption{Where the $z$ of an exact repeat lands if the original study had $p=0.05$
and the true effect equals the observed one. The bell is centred on the cut-off, so
the repeat is a coin flip: half the time (orange) it comes back ``not significant''.}
\label{fig:ht-replicate}
\end{figure}

The cleanest way to understand what a p-value \emph{is} comes from asking how it
behaves when the null is true. Simulate $2000$ studies of $n=20$ Normal observations
where $H_0\colon\mu=0$ holds exactly, run a t-test in each, and plot the p-values.
Then do the same where the true mean is $0.5\sigma$.

\begin{figure}[H]
\centering
\begin{tikzpicture}
\begin{axis}[
  name=nul, width=6.3cm, height=5cm, ybar interval, xmin=0, xmax=1, ymin=0, ymax=150,
  xlabel={p-value}, ylabel={count (of 2000)}, title={$H_0$ true}, title style={font=\small},
  axis lines=left, xtick={0,0.5,1}, x tick label as interval=false, xmajorgrids=false,
  xticklabel style={/pgf/number format/fixed}, clip=false,
]
\addplot[fill=histfill, draw=distblue] coordinates {
(0,87)(0.05,93)(0.10,96)(0.15,99)(0.20,101)(0.25,117)(0.30,96)(0.35,106)(0.40,94)(0.45,96)
(0.50,82)(0.55,86)(0.60,120)(0.65,103)(0.70,119)(0.75,100)(0.80,107)(0.85,117)(0.90,83)(0.95,98)(1.0,98)};
\filldraw[fill=trapred, draw=trapred!70!black] (axis cs:0,0) rectangle (axis cs:0.05,87);
\draw[densityred, dashed, thick] (axis cs:0,100) -- (axis cs:1,100);
\node[densityred, font=\footnotesize, anchor=south east] at (axis cs:1,122) {flat line: 100 per bin};
\draw[trapred] (axis cs:0.025,60) -- (axis cs:0.2,30);
\node[trapred, font=\footnotesize, anchor=west, fill=white, inner sep=1pt] at (axis cs:0.2,30) {$p<0.05$: 87};
\end{axis}
\begin{axis}[
  at={(nul.east)}, anchor=west, xshift=1.2cm,
  width=6.3cm, height=5cm, ybar interval, xmin=0, xmax=1, ymin=0, ymax=1200,
  xlabel={p-value}, title={True mean $0.5\sigma$}, title style={font=\small},
  axis lines=left, xtick={0,0.5,1}, x tick label as interval=false, xmajorgrids=false,
  xticklabel style={/pgf/number format/fixed}, clip=false,
]
\addplot[fill=targetgreen!50, draw=targetgreen!60!black] coordinates {
(0,1102)(0.05,282)(0.10,146)(0.15,103)(0.20,69)(0.25,60)(0.30,35)(0.35,36)(0.40,30)(0.45,23)
(0.50,21)(0.55,17)(0.60,15)(0.65,17)(0.70,10)(0.75,8)(0.80,6)(0.85,11)(0.90,4)(0.95,5)(1.0,5)};
\filldraw[fill=trapred, draw=trapred!70!black] (axis cs:0,0) rectangle (axis cs:0.05,1102);
\node[trapred, font=\footnotesize, anchor=west, align=left] at (axis cs:0.08,900) {$p<0.05$: 1102 of 2000\\$=55\%$ $=$ the power};
\end{axis}
\end{tikzpicture}
\caption{Same test, two worlds; the red bar counts p-values below $0.05$. Left: when
$H_0$ is true, p-values are spread evenly over $[0,1]$ (about $100$ per bin of width
$0.05$; dashed line), so in theory exactly the fraction $\alpha$ fall below $\alpha$;
here $87$ of $2000$ did. Right: when there is a real effect, p-values pile up near
$0$; the fraction below $0.05$ ($55\%$ here) is the power.}
\label{fig:ht-uniform}
\end{figure}

\begin{whybox}
Take a continuous statistic $T$ with null CDF $F$, and a right-tailed p-value
$p = 1-F(T)$. Under $H_0$, $F(T)$ is uniform on $[0,1]$ (the probability integral
transform: a variable plugged into its own CDF is uniform). So for any $u$,
\[
  \P_{H_0}(p \le u) = \P\bigl(F(T)\ge 1-u\bigr) = u .
\]
This is what makes $\alpha$ meaningful: ``reject when $p\le 0.05$'' raises false
alarms in exactly $5\%$ of true-null studies. It also shows why p-values from true
nulls are pure noise. A null study is as likely to give $p=0.04$ as $p=0.94$.
\end{whybox}

This uniformity is the root of the multiple-testing problem in Chapter 11: run
$20$ tests of true nulls and you expect one $p<0.05$ by construction.

\section{Assumptions and robustness of the t-test}
\label{sec:ht-robust}

Every t-test assumes:
\begin{enumerate}
\item \textbf{Independence.} Observations do not influence each other (within and,
for two-sample tests, between groups). This is the assumption that matters most and
cannot be fixed by a bigger sample. Repeated measurements of the same user, or
requests sharing a cache, violate it and make $\SE$ too small.
\item \textbf{Approximate Normality of the mean.} Exact for Normal data; otherwise
supplied by the CLT once $n$ is large enough. How large depends on skewness.
\item \textbf{Equal variances}, for the pooled test only. Welch drops this.
\end{enumerate}

How robust is the test to non-Normal data? Simulate it. Draw $40{,}000$ samples
from a distribution whose true mean equals $\mu_0$, run a two-sided one-sample t-test
at $\alpha=0.05$, and record how often it falsely rejects. A perfect test rejects
$5\%$ of the time.

\begin{figure}[H]
\centering
\begin{tikzpicture}
\begin{axis}[
  width=11.5cm, height=6cm, xmode=log, log basis x=2, xmin=4, xmax=200,
  xtick={5,10,20,40,80,160}, xticklabels={5,10,20,40,80,160},
  ymin=0, ymax=20, xlabel={sample size $n$}, ylabel={actual Type I error (\%)},
  axis lines=left, legend style={font=\small}, legend pos=north east,
  grid=major, grid style={gray!20},
]
\addplot[thick, mark=*, distblue] coordinates {(5,5.2)(10,5.0)(20,5.1)(40,4.8)(80,5.2)(160,5.1)};
\addlegendentry{Normal data}
\addplot[thick, mark=square*, quotientorange] coordinates {(5,11.9)(10,10.0)(20,7.9)(40,6.8)(80,5.8)(160,5.4)};
\addlegendentry{Exponential (skewed)}
\addplot[thick, mark=triangle*, trapred] coordinates {(5,17.9)(10,16.1)(20,13.1)(40,10.9)(80,9.0)(160,7.5)};
\addlegendentry{Log-normal (very skewed)}
\draw[gray, dashed] (axis cs:4,5) -- (axis cs:200,5);
\node[gray!70!black, font=\footnotesize, anchor=north east] at (axis cs:200,4.6) {target: $5\%$};
\end{axis}
\end{tikzpicture}
\caption{False-alarm rate of a nominal $5\%$ one-sample t-test, from simulation (the
dashed line is what it should be). For Normal data it is right at every $n$. For
skewed data it is too high at small $n$ and slowly converges to $5\%$ as the CLT
takes hold. Strong skew (latency-like data) needs well over $100$ observations.}
\label{fig:ht-robust}
\end{figure}

The lesson for latency data, which are typically right-skewed with a long tail: a
t-test on $10$ requests can be misleading, while a t-test on thousands is fine for
the \emph{mean}. If the median or p99 is the real question, use the bootstrap
(Chapter 8) or a permutation test on that statistic instead.

Unequal variances combined with unequal group sizes are a sharper danger for the
pooled test. Again by simulation ($H_0$ true, $40{,}000$ repetitions, $\alpha=0.05$):

\begin{center}
\begin{tabular}{cccc}
\toprule
group sizes & SDs & pooled rejects & Welch rejects\\
\midrule
$25,\ 25$ & $1,\ 1$ & $5.0\%$ & $5.0\%$\\
$25,\ 25$ & $4,\ 1$ & $5.6\%$ & $5.2\%$\\
$10,\ 40$ & $4,\ 1$ & $29.8\%$ & $5.1\%$\\
$40,\ 10$ & $4,\ 1$ & $0.1\%$ & $5.1\%$\\
\bottomrule
\end{tabular}
\end{center}

When the small group has the large variance, pooling underestimates the SE (the big,
quiet group dominates $s_p$) and the test rejects six times too often. When the large
group has the large variance, the test becomes uselessly conservative. Welch stays at
$5\%$ throughout.

\begin{figure}[H]
\centering
\begin{tikzpicture}[
  q/.style={draw, rounded corners, fill=structurepurple!10, align=center, font=\small, text width=3.1cm, minimum height=0.9cm},
  a/.style={draw, rounded corners, fill=targetgreen!12, align=center, font=\small, text width=2.6cm, minimum height=0.9cm},
  >={Stealth}, every edge quotes/.style={font=\footnotesize, auto}]
\node[q] (q1) at (0,0) {Same units measured twice?};
\node[a] (pair) at (4.3,0) {Paired t-test on the differences};
\node[q] (q2) at (0,-1.7) {Comparing to a fixed value $\mu_0$?};
\node[a] (one) at (4.3,-1.7) {One-sample t-test};
\node[a] (welch) at (4.3,-3.4) {Welch two-sample t-test};
\node[q] (q3) at (8.6,-1.7) {Small $n$ and heavy skew, or not the mean?};
\node[a] (perm) at (8.6,-3.7) {Permutation test or bootstrap};
\draw[->] (q1) edge["yes"] (pair);
\draw[->] (q1) edge["no"] (q2);
\draw[->] (q2) edge["yes"] (one);
\draw[->] (q2) edge["no"] (welch);
\draw[->, dashed] (one) -- (q3);
\draw[->, dashed] (pair.east) -- (q3.north west);
\draw[->, dashed] (welch.east) -- (q3.south west);
\draw[->] (q3) edge["yes"] (perm);
\end{tikzpicture}
\caption{Choosing a test for a mean. Start top left. Whichever t-test you land on,
check the robustness question on the right; if the answer is yes, switch to a
resampling method.}
\label{fig:ht-choose}
\end{figure}

\begin{lstlisting}[language=Python]
from scipy import stats
stats.ttest_1samp(x, 120)                   # t=3.65, p=0.0053
stats.ttest_ind(old, new)                   # pooled: t=3.28, p=0.0083
stats.ttest_ind(old, new, equal_var=False)  # Welch:  t=3.28, p=0.0099
stats.ttest_rel(a, b)                       # paired: t=8.52, p=0.00006
\end{lstlisting}

\begin{connectionbox}
Comparing two ML models on the same test set is a paired problem: both models see the
same examples, so compare per-example losses (or per-fold scores) with a paired test
or a paired bootstrap. Treating the two sets of scores as independent throws away the
pairing and can hide a real improvement, exactly as in Section~\ref{sec:ht-paired}.
\end{connectionbox}

Chapter 10 extends these ideas from means to proportions, several groups and
categorical tables; Chapter 11 deals with effect sizes and what happens when many
tests are run at once.

\section{Summary}

\begin{center}
\small
\begin{tabular}{p{3.4cm}p{9.6cm}}
\toprule
Idea & In one line\\
\midrule
Null hypothesis & The ``nothing is going on'' world; chosen because it is specific enough to compute a distribution.\\
Null distribution & What the test statistic looks like in that world (by simulation, shuffling, or formula).\\
p-value & Fraction of the null world at least as extreme as the observed result.\\
$\alpha$ & False-alarm rate fixed in advance; reject when $p\le\alpha$.\\
Fail to reject & Absence of evidence, not evidence of absence.\\
Permutation test & Shuffle labels to build the null world; needs only randomization.\\
z-test / t-test & $(\text{estimate} - \text{null value})/\SE$, compared to $\Normal(0,1)$ or $t_{\text{df}}$.\\
Two-sample & Welch by default: $\SE=\sqrt{s_1^2/n_1+s_2^2/n_2}$.\\
Paired & One-sample test on differences; removes between-unit variance.\\
Power & $1-\beta\approx\Phi(\delta/\SE - z_{\alpha/2})$; grows with effect size, $n$, $\alpha$.\\
A/B sample size & $n=(z_{\alpha/2}+z_\beta)^2[p_1q_1+p_2q_2]/\delta^2$ per arm.\\
Duality & Reject $\mu_0$ at level $\alpha$ $\iff$ $\mu_0$ outside the $(1-\alpha)$ CI.\\
Under $H_0$ & p-values are uniform on $[0,1]$.\\
Robustness & Independence is essential; skew needs large $n$; unequal variances need Welch.\\
\bottomrule
\end{tabular}
\end{center}

\section{Exercises}

\begin{exercise}
IQ-style scores have $\sigma=15$. A sample of $n=36$ has $\xbar=104.5$. Test
$H_0\colon\mu=100$ two-sided and one-sided ($H_1\colon\mu>100$). What do you conclude at
$\alpha=0.05$ in each case?
\end{exercise}

\begin{exercise}
Group A $=\{3,5,7\}$, group B $=\{1,2,4\}$. List the $\binom{6}{3}$ relabellings and find
the exact one-sided permutation p-value for $\bar a-\bar b$. Why can this test never
give $p<0.05$?
\end{exercise}

\begin{exercise}
Five users rate an app before and after a redesign. The differences (after $-$ before)
are $2, 4, 1, 3, 5$. Run a paired t-test of $H_0\colon\mu_d=0$ at $\alpha=0.05$.
\end{exercise}

\begin{exercise}
You run $20$ independent tests at $\alpha=0.05$ and every null is true. What is the
probability of at least one ``significant'' result?
\end{exercise}

\begin{exercise}
In the A/B example ($10\%\to11\%$, $14{,}748$ per arm), by what factor does $n$ change
if (a) you want $90\%$ power instead of $80\%$; (b) you want to detect a lift of
$2$ points instead of $1$ (ignore the small change in $p(1-p)$)?
\end{exercise}

\begin{exercise}
A $95\%$ CI for a mean difference is $(2.1, 5.3)$. At $\alpha=0.05$, do you reject
$H_0\colon\delta=0$? $H_0\colon\delta=5$? Can you tell from this interval alone whether
$H_0\colon\delta=0$ is rejected at $\alpha=0.01$?
\end{exercise}

\begin{exercise}
A report says: ``$p=0.02$, so there is a $98\%$ chance the new feature increases
engagement.'' Rewrite the sentence correctly.
\end{exercise}

\begin{exercise}
In the robustness table, the pooled test rejects $29.8\%$ of the time with group
sizes $10,40$ and SDs $4,1$. Compute $s_p$ and the pooled SE using the true SDs, compare
with the Welch SE, and explain the inflation.
\end{exercise}

\subsection*{Answers}

\begin{enumerate}
\item $\SE=15/6=2.5$, $z=1.8$. Two-sided $p=0.072$: do not reject. One-sided
$p=0.036$: reject. This is why sides must be fixed in advance.
\item Over the $20$ labellings, the sum of group A takes the values
6, 7, 8, 8, 9, 9, 10, 10, 10, 11, 11, 12, 12, 12, 13, 13, 14, 14, 15, 16.
The observed sum is $15$; sums $\ge15$ occur $2$ times, so $p=2/20=0.10$. The smallest
possible p-value is $1/20=0.05$, reached only by the single most extreme labelling.
\item $\bar d=3$, $s_d=\sqrt{2.5}=1.581$, $\SE=0.707$, $t=4.24$ on $4$ df. The cut-off is
$2.776$, and $p=0.013$: reject.
\item $1-0.95^{20}=0.64$.
\item (a) $(1.960+1.282)^2/(1.960+0.842)^2=10.51/7.85\approx1.34$, so about $34\%$ more
users. (b) $n\propto1/\delta^2$, so a quarter as many: about $3{,}700$ per arm.
\item Reject $\delta=0$ ($0$ is outside). Do not reject $\delta=5$ (inside). For
$\alpha=0.01$ you would need the wider $99\%$ interval, which this one does not give;
you cannot tell for sure.
\item ``If the feature had no effect, data showing an increase this large or larger
would occur about $2\%$ of the time.'' The probability that the feature works
requires a prior and Bayes' rule (Chapter 14).
\item $s_p^2=(9\cdot16+39\cdot1)/48=3.81$, $s_p=1.95$, pooled
$\SE=1.95\sqrt{1/10+1/40}=0.690$. Welch $\SE=\sqrt{16/10+1/40}=1.275$. The pooled SE is
about half the true one, so $|t|$ is about twice as large as it should be and the test
rejects far too often.
\end{enumerate}
